% GENERATED FILE. Edit canonical lessons, then run npm run generate:books.
% canonical-source-hash: fnv1a64:20a8166215327c5a
% canonical-lessons: MR-C238-vaimanik, MR-C238-samsodhak, MR-C238-pardesi, MR-C238-paryatak, MR-C238-praudh

\chapter{Pilot, Researcher, Foreigner, Tourist, Adult}
\label{ch:mr-pilot-researcher-foreigner-tourist-adult}
\hlchaptermodality{238}

\begin{chapteropening}
\textbf{By the end of this chapter:} \emph{I can say a pilot, a researcher, a foreigner, a tourist and an adult.}

Complete the last lesson of chapter 238: I can say a pilot, a researcher, a foreigner, a tourist and an adult.
\end{chapteropening}

\section[\texorpdfstring{\mr{वैमानिक} (vaimānik)}{वैमानिक (vaimānik)}]{\mr{वैमानिक} (vaimānik) — a pilot}

\label{lesson:MR-C238-vaimanik}



\begin{quote}
Before the new one: say the Marathi for a website, then the Marathi for a keyboard.
\end{quote}

\subsection*{You'll want to know: \mr{वैमानिक}}
\textbf{\mr{वैमानिक}} — \emph{vaimānik} — \textquotedblleft{}a pilot\textquotedblright{}.

\begin{grammarlens}[title={What you've built}]
One more word for this chapter.
\end{grammarlens}

\subsection*{Your turn}
\begin{itemize}
\raggedright
  \item \emph{Say it:} \emph{vaimānik}
  \item \emph{Say it:} \emph{vaimānik}, once more
  \item \emph{Say it:} say \emph{vaimānik}
  \item \emph{Recall:} say the Marathi for a website, then the Marathi for a keyboard, then say \emph{vaimānik} again
\end{itemize}

\begin{tcolorbox}[breakable,colback=teal!4,colframe=teal!35!black,title={Before you move on}]
What does \mr{वैमानिक} mean? (\textquotedblleft{}A pilot\textquotedblright{}.) Say it once more.
\end{tcolorbox}
\section[\texorpdfstring{\mr{संशोधक} (saṁśodhak)}{संशोधक (saṁśodhak)}]{\mr{संशोधक} (saṁśodhak) — a researcher}

\label{lesson:MR-C238-samsodhak}



\begin{quote}
Before the new one: say the Marathi for a keyboard, then the Marathi for a pilot.
\end{quote}

\subsection*{You'll want to know: \mr{संशोधक}}
\textbf{\mr{संशोधक}} — \emph{saṁśodhak} — \textquotedblleft{}a researcher\textquotedblright{}.

\begin{grammarlens}[title={What you've built}]
One more word for this chapter.
\end{grammarlens}

\subsection*{Your turn}
\begin{itemize}
\raggedright
  \item \emph{Say it:} \emph{saṁśodhak}
  \item \emph{Say it:} \emph{saṁśodhak}, once more
  \item \emph{Say it:} say \emph{saṁśodhak}
  \item \emph{Recall:} say the Marathi for a keyboard, then the Marathi for a pilot, then say \emph{saṁśodhak} again
\end{itemize}

\begin{tcolorbox}[breakable,colback=teal!4,colframe=teal!35!black,title={Before you move on}]
What does \mr{संशोधक} mean? (\textquotedblleft{}A researcher\textquotedblright{}.) Say it once more.
\end{tcolorbox}
\section[\texorpdfstring{\mr{परदेशी} (pardeśī)}{परदेशी (pardeśī)}]{\mr{परदेशी} (pardeśī) — a foreigner}

\label{lesson:MR-C238-pardesi}



\begin{quote}
Before the new one: say the Marathi for a pilot, then the Marathi for a researcher.
\end{quote}

\subsection*{You'll want to know: \mr{परदेशी}}
\textbf{\mr{परदेशी}} — \emph{pardeśī} — \textquotedblleft{}a foreigner\textquotedblright{}.

\begin{grammarlens}[title={What you've built}]
One more word for this chapter.
\end{grammarlens}

\subsection*{Your turn}
\begin{itemize}
\raggedright
  \item \emph{Say it:} \emph{pardeśī}
  \item \emph{Say it:} \emph{pardeśī}, once more
  \item \emph{Say it:} say \emph{pardeśī}
  \item \emph{Recall:} say the Marathi for a pilot, then the Marathi for a researcher, then say \emph{pardeśī} again
\end{itemize}

\begin{tcolorbox}[breakable,colback=teal!4,colframe=teal!35!black,title={Before you move on}]
What does \mr{परदेशी} mean? (\textquotedblleft{}A foreigner\textquotedblright{}.) Say it once more.
\end{tcolorbox}
\section[\texorpdfstring{\mr{पर्यटक} (paryaṭak)}{पर्यटक (paryaṭak)}]{\mr{पर्यटक} (paryaṭak) — a tourist}

\label{lesson:MR-C238-paryatak}



\begin{quote}
Before the new one: say the Marathi for a researcher, then the Marathi for a foreigner.
\end{quote}

\subsection*{You'll want to know: \mr{पर्यटक}}
\textbf{\mr{पर्यटक}} — \emph{paryaṭak} — \textquotedblleft{}a tourist\textquotedblright{}.

\begin{grammarlens}[title={What you've built}]
One more word for this chapter.
\end{grammarlens}

\subsection*{Your turn}
\begin{itemize}
\raggedright
  \item \emph{Say it:} \emph{paryaṭak}
  \item \emph{Say it:} \emph{paryaṭak}, once more
  \item \emph{Say it:} say \emph{paryaṭak}
  \item \emph{Recall:} say the Marathi for a researcher, then the Marathi for a foreigner, then say \emph{paryaṭak} again
\end{itemize}

\begin{tcolorbox}[breakable,colback=teal!4,colframe=teal!35!black,title={Before you move on}]
What does \mr{पर्यटक} mean? (\textquotedblleft{}A tourist\textquotedblright{}.) Say it once more.
\end{tcolorbox}
\section[\texorpdfstring{\mr{प्रौढ} (prauḍh)}{प्रौढ (prauḍh)}]{\mr{प्रौढ} (prauḍh) — an adult}

\label{lesson:MR-C238-praudh}



\begin{quote}
Before the new one: say the Marathi for a foreigner, then the Marathi for a tourist.
\end{quote}

\subsection*{You'll want to know: \mr{प्रौढ}}
\textbf{\mr{प्रौढ}} — \emph{prauḍh} — \textquotedblleft{}an adult\textquotedblright{}.

\begin{grammarlens}[title={What you've built}]
One more word for this chapter.
\end{grammarlens}

\subsection*{Your turn}
\begin{itemize}
\raggedright
  \item \emph{Say it:} \emph{prauḍh}
  \item \emph{Say it:} \emph{prauḍh}, once more
  \item \emph{Say it:} say \emph{prauḍh}
  \item \emph{Recall:} say the Marathi for a foreigner, then the Marathi for a tourist, then say \emph{prauḍh} again
\end{itemize}

\begin{tcolorbox}[breakable,colback=teal!4,colframe=teal!35!black,title={Before you move on}]
What does \mr{प्रौढ} mean? (\textquotedblleft{}An adult\textquotedblright{}.) Say it once more.
\end{tcolorbox}
